\Needspace{14\baselineskip}
\subsection{Densities from a compact K\"ahler factor}
\label{subsec:closed-factor-product-densities}

The proof of \cref{thm:zero-lower-mass-product-equation} used only two features of the projective factor and of the density $q_\alpha$. The factor is K\"ahler, so that Stokes' theorem fixes the normalization scalars there, and the potential lives on the factor, so that a single term of the binomial expansion survives on the product. Neither the radial form of $q_\alpha$ nor any bound on it beyond integrability is needed. We therefore replace $\mathbb P^k$ by an arbitrary compact K\"ahler manifold and solve the threshold equation for every integrable density on it.

Let $(Y,\kappa)$ be a connected compact K\"ahler manifold of dimension $k\geq1$ with $\int_Y\kappa^k=1$, and keep $(X_0,h_0)$, $m$, $V_0$ and $C$ as above. On $X=X_0\times Y$, of dimension $N=m+k$, we use the background $H=h_0+\kappa$, a smooth positive probability volume form $\mathrm{d}V$, the smooth positive probability volume form $\mu=V_0^{-1}h_0^m\wedge\kappa^k$ and the ratio $R=\mu/\mathrm{d}V$, and we write $\pi_Y\colon X\to Y$ for the projection. In the expansion of $H^N$ only the term $Ch_0^m\wedge\kappa^k$ survives, so
\[
 \int_XH^N=CV_0.
\]
For $0\leq g\in L^1(Y,\kappa^k)$ we put $M_g=\int_Yg\kappa^k$ and identify $g$ with $\pi_Y^*g$. By Fubini's theorem, $\int_XRg\,\mathrm{d}V=\int_Xg\,\mu=M_g$. The hypothesis $v_-(h_0)=0$ made at the beginning of the section is used below only to show that $H$ is not closed and that $v_-(H)=0$.

\Needspace{8\baselineskip}
\begin{lemma}\label{lem:closed-factor-bounded-normalization}\index{Normalization scalar!Kahler factor densities@K\"ahler factor densities}
Let $g\geq0$ be a bounded measurable function on $Y$ with $M_g>0$. There is a continuous $v\in\PSH(Y,\kappa)$ with $(\kappa+\ddc v)^k=M_g^{-1}g\kappa^k$, and
\[
 (H+\ddc\pi_Y^*v)^N=\frac{CV_0}{M_g}Rg\,\mathrm{d}V,
 \qquad
 c_H(Rg)=\frac{CV_0}{M_g}.
\]
\end{lemma}
\begin{proof}
By \cref{thm:kn-continuous} on $Y$, with background $\kappa$ and volume form $\kappa^k$, there are $c>0$ and a continuous $v\in\PSH(Y,\kappa)$ with $(\kappa+\ddc v)^k=cg\kappa^k$. Since $\kappa$ is closed, Stokes' theorem gives every bounded $\kappa$-plurisubharmonic function the mass $\int_Y\kappa^k=1$ (for $k\geq2$ this is part of \cref{thm:mass-invariance-characterization}); hence $cM_g=1$. The pullback $\pi_Y^*v$ is continuous and $H$-plurisubharmonic, since $H+\ddc\pi_Y^*v=h_0+\pi_Y^*(\kappa+\ddc v)$. In the binomial expansion of its bounded product (\cref{prop:bounded-products}), a power of $h_0$ of degree larger than $m$ vanishes, and so does a power of $\pi_Y^*(\kappa+\ddc v)$ of degree larger than $k$, since this current depends only on the $k$-dimensional factor. As in Step~2 of the proof of \cref{thm:zero-lower-mass-product-equation}, only one term remains:
\[
 (H+\ddc\pi_Y^*v)^N=Ch_0^m\wedge\pi_Y^*(\kappa+\ddc v)^k
 =\frac{C}{M_g}g\,h_0^m\wedge\kappa^k=\frac{CV_0}{M_g}Rg\,\mathrm{d}V.
\]
Uniqueness of the scalar among continuous solutions (\cref{prop:vanishing-exponent}) gives $c_H(Rg)=CV_0/M_g$.
\end{proof}

The lemma contains \cref{lem:split-product-normalization}, and it allows $g$ to vanish on a set of positive measure. Monotone truncation now determines the threshold of every integrable density from the factor, and the closed factor turns a threshold subsolution into an exact solution.

\Needspace{12\baselineskip}
\begin{theorem}\label{thm:closed-factor-product-equation}\index{Monge--Ampere equation@Monge--Amp\`ere equation!Kahler factor densities@K\"ahler factor densities}\index{Subsolution!threshold solution!Kahler factor@K\"ahler factor}
Let $0\leq g\in L^1(Y,\kappa^k)$ be nonzero, and put $f=Rg$. Then
\[
 b_H(f)=\frac{CV_0}{M_g}=\frac{\int_XH^N}{\int_Xf\,\mathrm{d}V}>0.
\]
Every $v\in\PSH(Y,\kappa)$ with $\mathrm{MA}^{\mathrm{ext}}_\kappa(v)\geq M_g^{-1}g\kappa^k$ satisfies
\[
 \mathrm{MA}^{\mathrm{ext}}_\kappa(v)=M_g^{-1}g\kappa^k,
 \qquad
 \mathrm{MA}^{\mathrm{ext}}_H(\pi_Y^*v)=b_H(f)f\,\mathrm{d}V,
\]
and the second measure has total mass $CV_0$. Such a $v$ exists in $\mathcal E_{\mathrm{env}}(Y,\kappa,0)$ with $\sup_Yv=0$; for it, $u=\pi_Y^*v$ belongs to $\mathcal E_{\mathrm{env}}(X,H,0)$, and $P_{\delta H}(u)\not\equiv-\infty$ for every $0<\delta<1$. Moreover, $H$ is not closed and $v_-(H)=0$.
\end{theorem}
\begin{proof}
\emph{Step 1: the threshold.} Write $a=\min_XR$, $A=\max_XR$ and $M_{g,s}=\int_Y\min(g,s)\kappa^k$ for $s>0$, so that $M_{g,s}\uparrow M_g$ as $s\to\infty$ by monotone convergence. Since $a\leq R\leq A$, the truncations of $f$ satisfy
\[
 R\min(g,j/A)\leq\min(f,j)\leq R\min(g,j/a).
\]
The functions $\min(g,s)$ are bounded and nonzero, so \cref{lem:closed-factor-bounded-normalization} and the monotonicity of the scalar in \cref{lem:bounded-normalizer-subsolution} give
\[
 \frac{CV_0}{M_{g,j/a}}\leq c_H(\min(f,j))\leq\frac{CV_0}{M_{g,j/A}}.
\]
Letting $j\to\infty$ gives $b_H(f)=CV_0/M_g$, and the second expression follows from $\int_XH^N=CV_0$ and $\int_Xf\,\mathrm{d}V=M_g$.

\emph{Step 2: the equation on the factor.} We use the results of \cref{sec:individual-density-subsolutions} on $Y$, with background $\kappa$ and volume form $\kappa^k$. Since $\kappa$ is K\"ahler, \cref{prop:individual-normalizer-lp} gives $b_\kappa(g)=1/M_g$, and by \cref{thm:individual-density-subsolution} some $v\in\PSH(Y,\kappa)$ satisfies $\mathrm{MA}^{\mathrm{ext}}_\kappa(v)\geq M_g^{-1}g\kappa^k$. Let $v$ be any such function. Its extended measure has total mass at most one: each bounded truncation $\max(v,-t)$ has mass $\int_Y\kappa^k=1$, restricting its Monge--Amp\`ere measure to $\{v>-t\}$ can only decrease this mass, and by \eqref{eq:extended-ma} these restrictions increase to $\mathrm{MA}^{\mathrm{ext}}_\kappa(v)$. Hence $\mathrm{MA}^{\mathrm{ext}}_\kappa(v)-M_g^{-1}g\kappa^k$ is a positive measure of total mass at most $1-1=0$, so it vanishes. Moreover, $v_-(\kappa)=\int_Y\kappa^k>0$, so \cref{cor:positive-lower-mass-exact-threshold}, applied on $Y$, gives such a $v$ in $\mathcal E_{\mathrm{env}}(Y,\kappa,0)$. Adding a constant, which keeps it in this class by \cref{prop:relative-elementary}, we may assume $\sup_Yv=0$.

\emph{Step 3: the equation on the product.} Let $v$ be as in Step~2 and $u=\pi_Y^*v$. Then $u\in\PSH(X,H)$, because $H+\ddc u=h_0+\pi_Y^*(\kappa+\ddc v)$. For $t>0$ we have $\max(u,-t)=\pi_Y^*\max(v,-t)$ and $\{u>-t\}=\pi_Y^{-1}(\{v>-t\})$, so the expansion in the proof of \cref{lem:closed-factor-bounded-normalization} gives
\[
 \mathbf1_{\{u>-t\}}\mathrm{MA}_H(\max(u,-t))
 =Ch_0^m\wedge\pi_Y^*\bigl(\mathbf1_{\{v>-t\}}\mathrm{MA}_\kappa(\max(v,-t))\bigr).
\]
Both sides increase with $t$, and \eqref{eq:extended-ma} on $X$ and on $Y$ gives
\[
 \mathrm{MA}^{\mathrm{ext}}_H(u)=Ch_0^m\wedge\pi_Y^*\mathrm{MA}^{\mathrm{ext}}_\kappa(v)
 =\frac{C}{M_g}g\,h_0^m\wedge\kappa^k=\frac{CV_0}{M_g}f\,\mathrm{d}V
 =b_H(f)f\,\mathrm{d}V.
\]
Its total mass is $b_H(f)\int_Xf\,\mathrm{d}V=CV_0$.

\emph{Step 4: the envelopes.} Suppose that $v\in\mathcal E_{\mathrm{env}}(Y,\kappa,0)$ and $\sup_Yv=0$. For $b\geq1$, \cref{def:relative-envelope-class} gives $w_b\in\PSH(Y,\kappa)$ with $b^{-1}w_b\leq v$. Since $H+\ddc\pi_Y^*w_b=h_0+\pi_Y^*(\kappa+\ddc w_b)\geq0$, the pullback of $w_b$ is $H$-plurisubharmonic, and $b^{-1}\pi_Y^*w_b\leq u$. As $u\leq0$, this shows that $u\in\mathcal E_{\mathrm{env}}(X,H,0)$. For $0<\delta<1$ take $b=1/\delta$. The function $\delta\pi_Y^*w_b$ is $\delta H$-plurisubharmonic, because $\delta H+\ddc(\delta\pi_Y^*w_b)=\delta h_0+\delta\pi_Y^*(\kappa+\ddc w_b)$, and it lies below $u$. Hence $P_{\delta H}(u)\geq\delta\pi_Y^*w_b$ is not identically $-\infty$.

Finally, Step~3 of the proof of \cref{thm:zero-lower-mass-product-equation} does not use the projective factor; it shows that $H$ is not closed and that $v_-(H)=0$.
\end{proof}

The threshold is thus given by the same expression as for a K\"ahler background in \cref{prop:individual-normalizer-lp}, although $H$ is not closed and has zero lower mass, and it depends only on the integral of the density. The measure $f\,\mathrm{d}V=g\mu$ is the product of the fixed volume form $V_0^{-1}h_0^m$ on $X_0$ with the measure $g\kappa^k$ on $Y$, so its dependence on $X_0$ is smooth and fixed, whereas the masses of bounded potentials collapse on $X_0$. Densities that also vary along $X_0$ behave differently: those constructed in \cref{thm:zero-threshold-solvable-concentration} have zero threshold.

For $Y=\mathbb P^k$ and $g=q_\alpha$, the theorem recovers the threshold and the equation of \cref{thm:zero-lower-mass-product-equation}. If $X_0$ is the elliptic threefold with the balanced metric of \cref{cor:elliptic-fivefold-exact-threshold}, and $(Y,\kappa)$ is any connected compact K\"ahler surface with $\int_Y\kappa^2=1$, then every probability density $g$ on $Y$ satisfies $b_H(Rg)=240\pi$ and has an exact threshold equation of mass $240\pi$. The theorem asserts nothing about uniqueness of the potentials on $X$, and nothing about densities that depend on both factors.
